% starter-ar.tex — Arabic edition of ch. 5, making a sourdough starter
%
% Original: The Sourdough Framework by Hendrik Kleinwächter,
% the-sourdough-framework.com, CC BY-SA 4.0. The Arabic edition inherits
% CC BY-SA 4.0.
%
% Arabic layer of upstream sourdough-starter/sourdough-starter.tex
% (scaffold structure, \label{ch:sourdough-starter},
% \label{sec:bakers-math}, figure/flowchart labels, the four \cite keys,
% footnote positions and the quoting epigraph verbatim; text re-typeset
% from the recovery reference text/05-sourdough-starter.txt, Phase 28R
% unit 7 — no retranslation). \qty{}/\qtyrange{}/\numproduct{} ride \LR{}
% islands (the ch. 7 \pHvalue precedent); bare \num{} digit runs stand
% per the basics-ar convention. \ch{CO2} rides \LR{} (history-ar
% precedent). Latin terms ride \textenglish{} islands.
%
% Cross-refs \ref{ch:wheat-sourdough} and \ref{ch:non-wheat-sourdough}
% are unresolved in this unit build (their chapters arrive in later
% units); they resolve at R3 assembly. Same pending set as earlier
% units.
%
% Tables: % table-bakers-math-example-ar.tex — Arabic edition of upstream
% tables/table-bakers-math-example.tex (ch. 5, جدول 5.1)
%
% Original: The Sourdough Framework by Hendrik Kleinwächter,
% the-sourdough-framework.com, CC BY-SA 4.0. The Arabic edition inherits
% CC BY-SA 4.0.
%
% RTL table mirror per the standing convention: the column order is the
% REVERSE of upstream (label column last), so the labels sit rightmost as
% in the accepted edition's render of p. 34 (book-ar-fixed.pdf,
% b78c2a28…): visual order left→right is percentage / هو / calculation /
% percentage / amount / label. Header cells keep the upstream \thead
% constructs (their centering behavior is identical either way); every
% \qty cell rides an \LR island. Cell contents recovered from the
% accepted render: header المكوّن / النسبة / الحساب, calculation cells
% «1000 g من 1000 g», the copula «هو», and the percentage column —
% label recovery, not retranslation. Geometry verified per-char against
% the accepted page after the 2026-11-15 mirror fix.

\begin{tabular}{@{}rcrrrl@{}}
\toprule
             &  & \thead{الحساب} & \thead{النسبة} &  & \thead{المكوّن} \\ \midrule
\LR{\qty{100}{\percent}} & هو & \LR{\qty{1000}{g}} من \LR{\qty{1000}{g}} & \LR{\qty{100}{\percent}} & \LR{\qty{1000}{g}} & دقيق \\
\LR{\qty{60}{\percent}}  & هو & \LR{\qty{600}{g}} من \LR{\qty{1000}{g}}  & \LR{\qty{60}{\percent}}  & \LR{\qty{600}{g}}  & ماء \\
\LR{\qty{10}{\percent}}  & هو & \LR{\qty{100}{g}} من \LR{\qty{1000}{g}}  & \LR{\qty{10}{\percent}}  & \LR{\qty{100}{g}}  & البادئة \\
\LR{\qty{2}{\percent}}   & هو & \LR{\qty{20}{g}} من \LR{\qty{1000}{g}}   & \LR{\qty{2}{\percent}}   & \LR{\qty{20}{g}}   & ملح \\ \bottomrule
\end{tabular}
 and
% % table-recipe-bakers-math-ar.tex — Arabic edition of upstream
% tables/table-recipe-bakers-math.tex (ch. 5, جدول 5.2)
%
% Original: The Sourdough Framework by Hendrik Kleinwächter,
% the-sourdough-framework.com, CC BY-SA 4.0. The Arabic edition inherits
% CC BY-SA 4.0.
%
% RTL table mirror per the standing convention, as
% table-bakers-math-example-ar.tex: the column order is the REVERSE of
% upstream (label column last, labels rightmost), matching the accepted
% edition's render of p. 34 (book-ar-fixed.pdf, b78c2a28…): visual order
% left→right is result / = / factor / × / 1400 / percentage / label. The
% \multicolumn{5}{l} header span is kept as upstream (it centers over its
% span under the TRT cell patch, as the accepted render shows); each cell
% holds a single run, so the \LR islands cover the quantities and the
% bare \times / = cells stay single-symbol. Cell contents recovered from
% the accepted render: header المكوّن / الحساب الخبازي / القيمة المحسوبة
% and the calculation columns 1400 × factor = result — label recovery,
% not retranslation. Geometry verified per-char against the accepted page
% after the 2026-11-15 mirror fix.

\begin{tabular}{@{}rcrcrrl@{}}
\toprule
\multicolumn{5}{l}{\thead{القيمة المحسوبة}} & \thead{الحساب الخبازي} & \thead{المكوّن} \\ \midrule
\LR{\qty{1400}{g}} &=& \LR{1} &\times& \LR{1400} & \LR{\qty{100}{\percent}} & دقيق \\
\LR{\qty{840}{g}}  &=& \LR{0.6} &\times& \LR{1400} & \LR{\qty{60}{\percent}}  & ماء \\
\LR{\qty{140}{g}}  &=& \LR{0.1} &\times& \LR{1400} & \LR{\qty{10}{\percent}}  & البادئة \\
\LR{\qty{28}{g}}   &=& \LR{0.02} &\times& \LR{1400} & \LR{\qty{2}{\percent}}   & ملح \\ \bottomrule
\end{tabular}
 (RTL mirrors).
% Flowcharts: % fig-starter-process-ar.tex — Arabic edition of the sourdough starter
% process flowchart (ch. 5, المخطط الانسيابي 5.1)
%
% Original: The Sourdough Framework by Hendrik Kleinwächter,
% the-sourdough-framework.com, CC BY-SA 4.0. The Arabic edition inherits
% CC BY-SA 4.0.
%
% R1-FIGURES route: re-derived from the upstream post-#492 source
% figures/fig-starter-process.tex (node styles from
% figures/flowcharts_tikz.tex, which sourdough-ar.sty inputs, and the
% node/edge geometry verbatim), with the label text recovered from the
% accepted edition's render of p. 35 (book-ar-fixed.pdf, b78c2a28…) —
% label recovery against a visual reference, not retranslation. Line
% splits inside nodes follow the accepted edition's render exactly.
% Arabic label text is wrapped in \textarabic{}; digit/Latin runs ride
% \LR{} islands (see fig-sourdough-history-timeline-ar.tex for the
% probe evidence and the Ruling 1 uniform-Pagella closure, ACCEPTED
% 2026-11-15). The U+0611 ARABIC SIGN ALAYHE ASSALLAM in تخؑلّص is
% genuinely present in the accepted edition's text layer (verified
% per-character with coordinates and against the 300-dpi render) and is
% reproduced verbatim.

\begin{tikzpicture}[node distance = 3.5cm, auto]
  \node [start] (start_n) {\textarabic{امزج \LR{\qty{50}{\gram}} دقيق}\\
                           \textarabic{مع \LR{\qty{50}{\gram}} ماء،}\\
                           \textarabic{وقلّب جيدًا}};
  \node [block, right of=start_n, node distance = 3.5cm] (wait_n) {\textarabic{انتظر}\\ \LR{\qty{24}{\hour}}};
  \node [decision, right of=wait_n, node distance = 3.5cm] (readycheck_n) {\textarabic{جهوزية؟}};
  \node [block, below of=wait_n, node distance = 3.2cm] (feed_n) {\textarabic{غّذِ الخليط}};
  \node [decision, right of=feed_n, node distance = 3.5cm] (limitcheck_n) {\textarabic{\LR{>10} تغذيات؟}};
  \node [fail, right of=limitcheck_n, node distance = 3.5cm] (abort_n) {\textarabic{تخؑلّص من الكل.}\\ \textarabic{وابدأ من جديد.}};
  \node [success, right of=readycheck_n, node distance = 3.5cm] (final_n) {\textarabic{جاهزة للاستخدام}};

  \draw [line] (start_n) -- (wait_n);
  \draw [line] (wait_n) -- (readycheck_n);
  \draw [line] (feed_n) -- (wait_n);
  \draw [line] (readycheck_n) -- node {\textarabic{لا}} (limitcheck_n);
  \draw [line] (limitcheck_n) -- node (feedok_n) {\textarabic{لا}} (feed_n) ;
  \draw [line] (limitcheck_n) -- node {\textarabic{نعم}} (abort_n);
  \draw [line] (readycheck_n) -- node {\textarabic{نعم}} (final_n);

  % Legend: each line is ONE \textarabic{} run with the bold label inline,
  % align=right so the lines hang from the right edge as in the accepted
  % render. Separate \textbf{\textarabic{}} labels order LTR against the
  % following run inside the node's LTR paragraph and land at line END;
  % the accepted render has them at line START (rightmost), which the
  % single-run form reproduces. Residual route deviations (documented,
  % same family as Ruling 1): node text rides the scaled Amiri (10.79 pt,
  % as in the ACCEPTED unit 5 flowcharts) where the accepted edition's
  % standalone-figure builds show 10.04 pt, and the legend box hugs its
  % text where the accepted box carries ~61 pt of dead space at the left
  % (a standalone-route artifact with no source-side equivalent).
  \node [below of=feedok_n, node distance=2cm, align=right] (details2) [shape=rectangle, draw, fill=maingray]{%
    \textarabic{\textbf{جهوزية:} زاد الخليط حجمًا، وبه فقاعات، ورائحته حمض اللبنيك أو حمض الخليك.}\\
    \textarabic{\textbf{تغذية الخليط:} تخؑلّص من كل شيء عدا \LR{10 g}، وامزج فيه \LR{50 g} دقيق و\LR{50 g} ماء.}
  };
\end{tikzpicture}
,
% fig-starter-readiness-ar.tex, fig-starter-maintenance-ar.tex
% (R1-FIGURES route, label recovery from the accepted edition's renders
% of pp. 35/40/42). Photos: raster carry-overs, upstream widths verbatim
% (0.85 / 0.6 \textwidth).
%
% Short captions: all seven floats carry upstream short forms,
% translated; flagged for the R3 LOF reconciliation like earlier units.

\chapter{صنع بادئة الخبز الحامض}%
\label{ch:sourdough-starter}
\begin{quoting}
ستتعلم في هذا الفصل كيف تصنع بادئة الخبز الحامض بنفسك، لكن قبل ذلك ستتعرّف سريعًا على الحساب الخبازي. لا تقلق، فهو أسلوب بسيط جدًا في كتابة الوصفات، أنظف وأسهل في التوسيع. وما إن تعتاد عليه حتى ترغب في كتابة كل وصفة بهذه الطريقة. وستتعلم فهم المؤشرات الدالة على جهوزية بادئتك، وكيف تجهّز بادئتك للتخزين طويل الأمد.
\end{quoting}

\section{الحساب الخبازي}%
\label{sec:bakers-math}

في مخبزٍ كبير، يكون العامل الحاسم هو مقدار الدقيق المتوفر لديك. وبناءً على كمية الدقيق المتاحة، يمكنك حساب عدد الأرغفة أو أرغفة الخبز الصغير التي يمكنك صنعها. ولكي يسهل الأمر على الخبازين، تُحسب كمية كل مكوّن نسبةً مئويةً من كمية الدقيق المتاحة. دعني أوضح ذلك بمثال صغير من مطعم بيتزا. في الصباح تدقق في المخزون فتكتشف أن لديك نحو \LR{\qty{1}{\kg}} من الدقيق. وصفتك المعتادة تتطلب نحو \LR{\qty{600}{\gram}} من الماء؛ وسيكون ذلك عجين بيتزا نموذجيًا، ليس جافًا أكثر من اللازم ولا طريًا أكثر من اللازم. ثم ستستخدم نحو \LR{\qty{20}{\gram}} من الملح ونحو \LR{\qty{100}{\gram}} من بادئة الخبز الحامض.\footnote{هذه هي وصفة عجين البيتزا المفضلة لدي. وفي نابولي تستخدم المطاعم الحديثة الخميرة الطازجة أو الجافة، غير أن البيتزا قديمًا كانت تُصنع بالخبز الحامض دوامًا.} وفي اليوم التالي تجد فجأة \LR{\qty{1.4}{\kg}} من الدقيق بين يديك، وبالتالي تستطيع صنع عجين بيتزا أكثر. فماذا تفعل؟ هل تضرب كل المكونات في \num{1.4}؟ نعم يمكن ذلك، لكن هناك طريقة أسهل. وهنا يأتي دور الحساب الخبازي. لننظر إلى الوصفة المعتادة بصيغة الحساب الخبازي، ثم نعدّلها بما يوافق كمية \LR{\qty{1.4}{\kg}} من الدقيق.

\begin{table}[!htb]
\centering
  % table-bakers-math-example-ar.tex — Arabic edition of upstream
% tables/table-bakers-math-example.tex (ch. 5, جدول 5.1)
%
% Original: The Sourdough Framework by Hendrik Kleinwächter,
% the-sourdough-framework.com, CC BY-SA 4.0. The Arabic edition inherits
% CC BY-SA 4.0.
%
% RTL table mirror per the standing convention: the column order is the
% REVERSE of upstream (label column last), so the labels sit rightmost as
% in the accepted edition's render of p. 34 (book-ar-fixed.pdf,
% b78c2a28…): visual order left→right is percentage / هو / calculation /
% percentage / amount / label. Header cells keep the upstream \thead
% constructs (their centering behavior is identical either way); every
% \qty cell rides an \LR island. Cell contents recovered from the
% accepted render: header المكوّن / النسبة / الحساب, calculation cells
% «1000 g من 1000 g», the copula «هو», and the percentage column —
% label recovery, not retranslation. Geometry verified per-char against
% the accepted page after the 2026-11-15 mirror fix.

\begin{tabular}{@{}rcrrrl@{}}
\toprule
             &  & \thead{الحساب} & \thead{النسبة} &  & \thead{المكوّن} \\ \midrule
\LR{\qty{100}{\percent}} & هو & \LR{\qty{1000}{g}} من \LR{\qty{1000}{g}} & \LR{\qty{100}{\percent}} & \LR{\qty{1000}{g}} & دقيق \\
\LR{\qty{60}{\percent}}  & هو & \LR{\qty{600}{g}} من \LR{\qty{1000}{g}}  & \LR{\qty{60}{\percent}}  & \LR{\qty{600}{g}}  & ماء \\
\LR{\qty{10}{\percent}}  & هو & \LR{\qty{100}{g}} من \LR{\qty{1000}{g}}  & \LR{\qty{10}{\percent}}  & \LR{\qty{100}{g}}  & البادئة \\
\LR{\qty{2}{\percent}}   & هو & \LR{\qty{20}{g}} من \LR{\qty{1000}{g}}   & \LR{\qty{2}{\percent}}   & \LR{\qty{20}{g}}   & ملح \\ \bottomrule
\end{tabular}

  \caption[مثال الحساب الخبازي]{جدول مثال يوضح كيفية الحساب الصحيح باستخدام الحساب الخبازي.}
\end{table}

لاحظ كيف حُسب كل مكوّن نسبةً مئويةً من الدقيق. فالنسبة \LR{\qty{100}{\percent}} هي خط الأساس، وتمثل الكمية المطلقة من الدقيق المتوفر لديك. وفي هذه الحالة هي \LR{\qty{1000}{\gram}}~(\LR{\qty{1}{\kg}}).

ولنعد الآن إلى مثالنا ونعدّل كمية الدقيق، إذ لدينا في اليوم التالي دقيق أكثر. وكما ذكرنا، لدينا في اليوم التالي \LR{\qty{1.4}{\kg}} بين يدينا (\LR{\qty{1400}{\gram}}).

\begin{table}[!htb]
    \centering
        % table-recipe-bakers-math-ar.tex — Arabic edition of upstream
% tables/table-recipe-bakers-math.tex (ch. 5, جدول 5.2)
%
% Original: The Sourdough Framework by Hendrik Kleinwächter,
% the-sourdough-framework.com, CC BY-SA 4.0. The Arabic edition inherits
% CC BY-SA 4.0.
%
% RTL table mirror per the standing convention, as
% table-bakers-math-example-ar.tex: the column order is the REVERSE of
% upstream (label column last, labels rightmost), matching the accepted
% edition's render of p. 34 (book-ar-fixed.pdf, b78c2a28…): visual order
% left→right is result / = / factor / × / 1400 / percentage / label. The
% \multicolumn{5}{l} header span is kept as upstream (it centers over its
% span under the TRT cell patch, as the accepted render shows); each cell
% holds a single run, so the \LR islands cover the quantities and the
% bare \times / = cells stay single-symbol. Cell contents recovered from
% the accepted render: header المكوّن / الحساب الخبازي / القيمة المحسوبة
% and the calculation columns 1400 × factor = result — label recovery,
% not retranslation. Geometry verified per-char against the accepted page
% after the 2026-11-15 mirror fix.

\begin{tabular}{@{}rcrcrrl@{}}
\toprule
\multicolumn{5}{l}{\thead{القيمة المحسوبة}} & \thead{الحساب الخبازي} & \thead{المكوّن} \\ \midrule
\LR{\qty{1400}{g}} &=& \LR{1} &\times& \LR{1400} & \LR{\qty{100}{\percent}} & دقيق \\
\LR{\qty{840}{g}}  &=& \LR{0.6} &\times& \LR{1400} & \LR{\qty{60}{\percent}}  & ماء \\
\LR{\qty{140}{g}}  &=& \LR{0.1} &\times& \LR{1400} & \LR{\qty{10}{\percent}}  & البادئة \\
\LR{\qty{28}{g}}   &=& \LR{0.02} &\times& \LR{1400} & \LR{\qty{2}{\percent}}   & ملح \\ \bottomrule
\end{tabular}

        \caption[مثال آخر للحساب الخبازي]{وصفة مثال تعتمد \LR{\qty{1400}{\gram}} خطَّ أساس لها، ثم تُحسب باستخدام الحساب الخبازي.}
\end{table}

نحسب لكل مكوّن نسبته المئوية اعتمادًا على الدقيق المتاح (\LR{\qty{1400}{\gram}}). ففي الماء مثلًا نحسب \LR{\qty{60}{\percent}} من \num{1400}. افتح آلة حاسبتك وأدخل \LR{\numproduct{1400 x 0.6}} وستحصل على القيمة الدقيقة بالغرامات التي ينبغي استخدامها. وفي اليوم الثاني يكون ذلك \LR{\qty{840}{\gram}}. أجرِ العملية نفسها على سائر المكونات وستعرف وصفتك.

لنفترض أنك أردت في اليوم التالي استخدام \LR{\qty{50}{\kg}} من الدقيق. ماذا ستفعل؟ ستعيد ببساطة حساب النسب المئوية مرة أخرى. تعجبني هذه الطريقة في كتابة الوصفات كثيرًا. تخيّل أنك أردت إعداد مكرونة، وأردت معرفة مقدار الصلصة التي ينبغي أن تصنعها. وبدلًا من إعداد وصفة لك وحدك، تصل فجأة عائلة جائعة. وأنت مكلف بإعداد مكرونة لِنحو \num{20} شخصًا. كيف ستحسب كمية الصلصة اللازمة؟ تذهب إلى الإنترنت وتفتش عن وصفة، ثم تضيع تمامًا في محاولة تكبيرها.
\section{عملية صنع البادئة}

\begin{figure}[!htb]
  \centering
  \includegraphics[width=0.85\textwidth]{sourdough-starter-activity-indicators}
  \caption[بادئة خبز حامض نشيطة جدًا]{بادئة خبز حامض نشيطة جدًا تظهرها الفقاعات في العجين.}%
  \label{fig:sourdough-starter}
\end{figure}

صنع بادئة الخبز الحامض أمر سهل جدًا؛ كل ما تحتاجه هو قليل من الصبر. بل هو سهل لدرجة يمكن تلخيصه في مخطط تدفق واحد~\ref{fig:sourdough-starter-process}. والدقيق المثالي لإنشاء بادئتك هو الدقيق الكامل. يمكنك استخدام دقيق القمح الكامل، أو الجاودار الكامل، أو السبلت الكاملة، أو أي دقيق كامل آخر لديك. بل إن الدقائق الخالية من الغلوتين، كدقيق الأرز أو الذرة، تصلح أيضًا. لا تقلق، فبإمكانك تغيير الدقيق لاحقًا متى شئت. استخدم أي دقيق كامل لديك أصلًا.

\begin{flowchart}[!htb]
\centering
  % fig-starter-process-ar.tex — Arabic edition of the sourdough starter
% process flowchart (ch. 5, المخطط الانسيابي 5.1)
%
% Original: The Sourdough Framework by Hendrik Kleinwächter,
% the-sourdough-framework.com, CC BY-SA 4.0. The Arabic edition inherits
% CC BY-SA 4.0.
%
% R1-FIGURES route: re-derived from the upstream post-#492 source
% figures/fig-starter-process.tex (node styles from
% figures/flowcharts_tikz.tex, which sourdough-ar.sty inputs, and the
% node/edge geometry verbatim), with the label text recovered from the
% accepted edition's render of p. 35 (book-ar-fixed.pdf, b78c2a28…) —
% label recovery against a visual reference, not retranslation. Line
% splits inside nodes follow the accepted edition's render exactly.
% Arabic label text is wrapped in \textarabic{}; digit/Latin runs ride
% \LR{} islands (see fig-sourdough-history-timeline-ar.tex for the
% probe evidence and the Ruling 1 uniform-Pagella closure, ACCEPTED
% 2026-11-15). The U+0611 ARABIC SIGN ALAYHE ASSALLAM in تخؑلّص is
% genuinely present in the accepted edition's text layer (verified
% per-character with coordinates and against the 300-dpi render) and is
% reproduced verbatim.

\begin{tikzpicture}[node distance = 3.5cm, auto]
  \node [start] (start_n) {\textarabic{امزج \LR{\qty{50}{\gram}} دقيق}\\
                           \textarabic{مع \LR{\qty{50}{\gram}} ماء،}\\
                           \textarabic{وقلّب جيدًا}};
  \node [block, right of=start_n, node distance = 3.5cm] (wait_n) {\textarabic{انتظر}\\ \LR{\qty{24}{\hour}}};
  \node [decision, right of=wait_n, node distance = 3.5cm] (readycheck_n) {\textarabic{جهوزية؟}};
  \node [block, below of=wait_n, node distance = 3.2cm] (feed_n) {\textarabic{غّذِ الخليط}};
  \node [decision, right of=feed_n, node distance = 3.5cm] (limitcheck_n) {\textarabic{\LR{>10} تغذيات؟}};
  \node [fail, right of=limitcheck_n, node distance = 3.5cm] (abort_n) {\textarabic{تخؑلّص من الكل.}\\ \textarabic{وابدأ من جديد.}};
  \node [success, right of=readycheck_n, node distance = 3.5cm] (final_n) {\textarabic{جاهزة للاستخدام}};

  \draw [line] (start_n) -- (wait_n);
  \draw [line] (wait_n) -- (readycheck_n);
  \draw [line] (feed_n) -- (wait_n);
  \draw [line] (readycheck_n) -- node {\textarabic{لا}} (limitcheck_n);
  \draw [line] (limitcheck_n) -- node (feedok_n) {\textarabic{لا}} (feed_n) ;
  \draw [line] (limitcheck_n) -- node {\textarabic{نعم}} (abort_n);
  \draw [line] (readycheck_n) -- node {\textarabic{نعم}} (final_n);

  % Legend: each line is ONE \textarabic{} run with the bold label inline,
  % align=right so the lines hang from the right edge as in the accepted
  % render. Separate \textbf{\textarabic{}} labels order LTR against the
  % following run inside the node's LTR paragraph and land at line END;
  % the accepted render has them at line START (rightmost), which the
  % single-run form reproduces. Residual route deviations (documented,
  % same family as Ruling 1): node text rides the scaled Amiri (10.79 pt,
  % as in the ACCEPTED unit 5 flowcharts) where the accepted edition's
  % standalone-figure builds show 10.04 pt, and the legend box hugs its
  % text where the accepted box carries ~61 pt of dead space at the left
  % (a standalone-route artifact with no source-side equivalent).
  \node [below of=feedok_n, node distance=2cm, align=right] (details2) [shape=rectangle, draw, fill=maingray]{%
    \textarabic{\textbf{جهوزية:} زاد الخليط حجمًا، وبه فقاعات، ورائحته حمض اللبنيك أو حمض الخليك.}\\
    \textarabic{\textbf{تغذية الخليط:} تخؑلّص من كل شيء عدا \LR{10 g}، وامزج فيه \LR{50 g} دقيق و\LR{50 g} ماء.}
  };
\end{tikzpicture}

  \caption[عملية صنع بادئة الخبز الحامض الكاملة]{عملية صنع بادئة الخبز الحامض من الصفر.}%
  \label{fig:sourdough-starter-process}
\end{flowchart}

دقيقك ملوّث بملايين الكائنات الدقيقة. وكما شرحنا سابقًا في فصل الخمائر والبكتيريا البرية، تعيش هذه الكائنات الدقيقة على سطح النبات. ولهذا يعمل الدقيق الكامل أفضل؛ إذ يكون التلوث الطبيعي بالكائنات الدقيقة التي تحاول تنميتها في بادئتك أكبر. فأكثر منها يعيش على القشرة مقارنةً بالكائنات الدقيقة الداخلية الساكنة في الحبة.

ابدأ بوزن نحو \LR{\qty{50}{\gram}} من الدقيق والقدر نفسه من الماء. لا يلزم أن تكون القياسات دقيقة؛ يمكنك استخدام أقل أو أكثر، أو تقدير النسب بالعين مجردة. فتلك القيم ليست سوى مرجع.

لا تستخدم ماءً معالجًا بالكلور عند إنشاء بادئتك. ومن الأفضل استخدام ماء معلَّب. وفي بعض المناطق، كألمانيا، يكون ماء الصنبور جيدًا تمامًا. يُضاف الكلور إلى الماء مطهرًا لقتل الكائنات الدقيقة؛ ولن تتمكن من تنمية بادئة بماء معالج بالكلور.

في هذه العملية، ترطيب بادئتك \LR{\qty{100}{\percent}}، أي أنك تستخدم كمية متساوية من الدقيق والماء. امزج كل شيء جيدًا حتى يترطب الدقيق كله ترطيبًا سليمًا. تفعّل هذه الخطوة الأبواغ الدقيقة في خليطك، فتخرجها من سباتها وتعيد إليها الحياة. وأخيرًا غطِّ الخليط مع التأكد أن الغطاء ليس محكم الإغلاق؛ فما زلت تريد إتاحة بعض تبادل الغازات. وأنا أحب استخدام كوب زجاجي وضعَ آخر مقلوب فوقه.

والآن تبدأ معركة ملحمية، كما يصوّرها الشكل~\ref{fig:sourdough-starter-microbial-war}. ففي إحدى الدراسات~\cite{yeasts+biocontrol+agent} حدد العلماء أكثر من \num{150} نوع مختلف من الخمائر يعيش على ورقة نبات واحدة. وكل الخمائر والبكتيريا المختلفة تحاول استمالة الأمر في هذه المعركة. كما تُنشَّط مسبّبات أمراض أخرى، كالعفن، بإضافة الماء. ولن ينجُ إلا أقوى الكائنات الدقيقة وأكثرها قابلية للتكيف.

\begin{figure}[!htb]
  \centering
  \includegraphics[width=0.6\textwidth]{sourdough-starter-microbial-war}
  \caption[حرب الكائنات الدقيقة في الأيام الأولى للخبز الحامض]{تصوّر مبسّط لحرب الكائنات الدقيقة التي تجري أثناء صنع بادئة الخبز الحامض. تنشط الأبواغ البرية الموجودة على النبات والدقيق لحظة خلط الدقيق بالماء. ولن ينجُ إلا كائنات الدقيق الأكثر تأقلمًا مع تخمير الدقيق. ونظرًا لعدم الرغبة في تخمير دقيق غير مرغوب فيه، يُنصح بالتخلص من البادئة الزائدة المتبقية عن التغذية في الأيام الأولى. وتحاول الخمائر والبكتيريا الناجية التفوق المستمر بعضها على بعض في المنافسة على الموارد. ويصعب على الكائنات الدقيقة الجديدة دخول البادئة، فتُقصى.}%
  \label{fig:sourdough-starter-microbial-war}
\end{figure}

فبإضافة الماء إلى الدقيق يبدأ النشاء بالتفكك. وتحاول الشتلة أن تنبت لكنها لم تعد قادرة. والأميلاز إنزيم أساسي في هذه العملية؛ إذ يُفكك النشاء المضغوط إلى سكريات أسهل هضمًا تغذي نمو النبات. والغلوكوز هو ما يحتاجه النبات لينمو. والكائنات الدقيقة التي تنجو من هذه الحماسة هي المتأقلمة مع استهلاك الغلوكوز.

وإلينا نحن الخبازين من الحظ أن الخمائر والبكتيريا تعرف جيدًا كيف تستقلب الغلوكوز؛ فهذا ما كانت النباتات تغذيها به في البرية. فقد كانت تشكّل بقعًا على الورقة وتحمي النبات من المسبّبات، في مقابل تلقيها الغلوكوز جزاءً على خدماتها. وتحاول كلُّ من الكائنات الدقيقة أن تهزم الأخرى: إما باستقلاب الغذاء بأسرع ما يمكن، أو بإنتاج عوامل تكبح استهلاك الغذاء من قبل غيرها، أو بإنتاج مبيدات بكتيرية و/أو فطرية. وهذه المرحلة المبكرة من البادئة بالغة الإثارة، إذ قد يكشف المزيد من البحث عن مبيدات فطرية أو مضادات حيوية جديدة.

وبحسب مصدر دقيقك، قد تختلف الكائنات الدقيقة المبدئية في بادئتك عن تلك الموجودة في بادئة أخرى. وقد أفاد بعض الناس أيضًا بأن كائنات يدك أو الهواء قد تؤثر في كائنات بادئتك الدقيقة. وهذا منطقي إلى حد ما. فكائنات يدك الدقيقة قد تكون بارعة في تخمير عرقك، لكنها على الأرجح ليست بارعة في استقلاب الغلوكوز. وقد يلعب تلوث يديك أو الهواء دورًا ثانويًا في المعركة الملحمية الأولية. لكن لن ينجُ إلا الكائنات الدقيقة الأكثر لياقة والملائمة لبيئة الخبز الحامض.

وهذا يعني أن الكائنات الدقيقة القادرة على تحويل المالتوز أو الغلوكوز ستحصد السبق. أو الكائنات التي تخمّر فضلات الكائنات الدقيقة الأخرى. فالإيثانول الذي تنتجه الخميرة تستقلبه البكتيريا في خبزك الحامض؛ ولهذا لا كحول في الخبز الحامض. ويمكنني تأكيد دور التلوث الجوي إلى حد ما؛ فحين أنشئ بادئة جديدة تكون العملية كلها سريعة بالنسبة لي، وبعد أيام قليلة تبدو بادئتي الجديدة حية بالفعل. وقد يكون ذلك بسبب تلوث مسبق في مطبخي بكائنات دقيقة مخمِّرة للدقيق.

انتظر نحو 24 ساعة وراقب ما يحدث لبادئتك. وقد ترى بعض العلامات المبكرة للتخمير بالفعل. استخدم أنفك لتشمّ العجين، وابحث عن فقاعات فيه؛ فقد زاد حجم عجينك قليلًا بالفعل. وما تراه وما تنتبه إليه علامةٌ على المعركة الأولى.

فبعض الكائنات الدقيقة تفوّق عليها غيرها بالفعل، وبعضها الآخر انتصر في المعركة الأولى. وبعد نحو \LR{\qty{24}{\hour}} يكون معظم النشاء قد تفكك، وتكون كائناتك الدقيقة جائعة إلى سكريات إضافية. فبالمغرفة خذ نحو \LR{\qty{10}{\gram}} من خليط اليوم السابق وضعها في وعاء جديد. ومرة أخرى، يمكنك ببساطة تقدير كل الكميات بالعين؛ فالأمر ليس بهذه الأهمية. امزج \LR{\qty{10}{\gram}} من اليوم السابق مع \LR{\qty{50}{\gram}} أخرى من الدقيق و\LR{\qty{50}{\gram}} من الماء.

لاحظ النسبة \LR{1:5}. فأنا كثيرًا ما أستخدم جزءًا واحدًا من الثقافة القديمة مع خمسة أجزاء من الدقيق وخمسة أجزاء من الماء. وهذه أيضًا هي النسبة التي أستخدمها غالبًا عند صنع العجين. فالعجين ليس سوى بادئة خبز حامض عملاقة بخصائص مختلفة قليلًا. وأنا دائمًا ما أستخدم نحو \LR{\qtyrange{100}{200}{\gram}} من البادئة لكل نحو \LR{\qty{1000}{\gram}} من الدقيق (بالحساب الخبازي: \LR{\qtyrange{10}{20}{\percent}}).

امزج خليطك الجديد مرة أخرى بالمغرفة حتى يتجانس. ثم غطِّه مرة أخرى بكوب أو بغطاء. وإن لاحظت أن سطح خليطك يجف كثيرًا ففكر في استخدام غطاء آخر؛ فالأجزاء الجافة سيتخلص منها كائنات دقيقة أكثر تأقلمًا، كالعفن. وفي تقارير مستخدمين كثيرة رأيت العفن قادرًا على الإضرار بالبادئة حين جفت البادئة نفسها جفافًا شديدًا.

سيبقى لديك بعض الخليط من يومك الأول. وحين يحتمل أن هذا الخليط يحوي مسبّبات أمراض خطرة نُشِّطت، تأكد من التخلص منه. والقاعدة العملية هي البدء في الاحتفاظ بالبادئة الزائدة لحظة صنع أول رغيف ناجح لك. فعندها تكون بادئتك الزائدة دقيقًا مخمَّرًا طويلًا، وهي إضافة ممتازة لصنع البسكويت المقرمش \textenglish{(crackers)} أو الأقراص المسطّحة أو خبز الشطائر اللذيذ الدسم… كما أجففها كثيرًا وأستخدمها عاملًا مضادًا للالتصاق عند فرد البيتزا التي أصنعها.\footnote{قد يكون التخلص من البادئة عند تحضير دفعة جديدة محبطًا. فمع الخبرة يصبح صنع الخبز أكثر كفاءة، ونادرًا ما تُنتَج بادئة زائدة. ومن الممكن تحضير كمية البادئة اللازمة لعجين الخبز بالضبط. بل إن البادئة المستنفدة كليًا يمكن إحياؤها بقطعة صغيرة من عجين الخبز. وأي بادئة زائدة متبقية، وهي غنية بالأبواغ، يمكن أن تكون احتياطًا لإنشاء بادئة جديدة. كل ما في الأمر أن تخلط البادئة الزائدة بقليل من الدقيق والماء، فتعود إلى الحياة. وهذا خيار عظيم إن استُنفدت البادئة سهوًا. ونهج عملي هو تخزين كل البادئة الزائدة في برطمان واحد في الثلاجة، وإضافة البادئة الجديدة فوقها كلما دعت الحاجة، واستخدامها وقت الحاجة.}

نأمل أن ترى من جديد بعض الفقاعات، وزيادة في حجم البادئة، أو تغيرًا في رائحتها. ويستسلم بعض الناس بعد اليوم الثاني أو الثالث، لأن العلامات قد لا تبقى بارزة كما كانت في اليوم الأول. والسبب في ذلك أن عددًا قليلًا من كائنات دقيقة مختارة بدأت تسيطر على البادئة بأكملها. فالأكثر تأقلمًا سيكون هو المنتصر، وهي بالغة الصغر في العدد وستنمو مع كل تغذية تالية. وحتى لو لم ترَ علامات نشاط مباشرة، فلا تقلق؛ فهناك نشاط في بادئتك على المستوى المجهري.

وبعد 24 ساعة أخرى نكرر الشيء نفسه مرة أخرى حتى نرى أن بادئة الخبز الحامض قد صارت نشيطة. والمزيد من ذلك في القسم التالي من هذا الكتاب.
\section{تحديد جهوزية البادئة}

عند بعض الناس يستغرق إنشاء البادئة كلها 4 أيام فقط. وعند آخرين قد تستغرق 7 أيام، وعند بعضهم حتى 20 يومًا. ويعتمد ذلك على عدة عوامل، من بينها مدى برعة كائناتك الدقيقة البرية في تخمير الدقيق. وبصورة عامة، تصبح بادئتك مع كل تغذية أكثر تأقلمًا مع بيئتها، وأفضل في تخمير الدقيق. ولهذا قد تكون البادئة العريقة الناضجة التي تتلقاها من صديق أقوى من بادئتك في البداية. ومع الوقت ستلحق بادئة الخبز الحامض عندك بالركب. وبالمثل، عُزلت خميرة الخَبز الحديثة هكذا من بادئات خبز حامض عمرها قرون.

\begin{flowchart}[!htb]
\centering
  % fig-starter-readiness-ar.tex — Arabic edition of the sourdough starter
% readiness flowchart (ch. 5, المخطط الانسيابي 5.2)
%
% Original: The Sourdough Framework by Hendrik Kleinwächter,
% the-sourdough-framework.com, CC BY-SA 4.0. The Arabic edition inherits
% CC BY-SA 4.0.
%
% R1-FIGURES route: re-derived from the upstream post-#492 source
% figures/fig-starter-readiness.tex (node styles from
% figures/flowcharts_tikz.tex and the node/edge geometry verbatim,
% including the decision_start→feed_yes_branch.north west and the
% size_check/smell_check south anchors), with the label text recovered
% from the accepted edition's render of p. 40 (book-ar-fixed.pdf,
% b78c2a28…) — label recovery against a visual reference, not
% retranslation. Line splits inside nodes follow the accepted edition's
% render exactly. The upstream r@{}l tabulars in the ratio nodes are set
% as plain \\ lines: the accepted edition's render shows four plain
% lines per ratio node. Arabic label text is wrapped in \textarabic{};
% digit/Latin runs ride \LR{} islands (Ruling 1 uniform-Pagella figure
% face, ACCEPTED 2026-11-15). The two feed-duration nodes render the
% range with a single trailing unit («6–12 h»), as the accepted render's
% per-char positions show; the caption keeps the repeated-unit
% \qtyrange form («6 h–12 h») that the accepted caption shows.

\begin{tikzpicture}[node distance = 3cm, auto]
  \node [start] (init) {\textarabic{أنشئ}\\ \textarabic{بادئة}};
  \node [decision_start, right of=init, node distance=3.5cm] (decision_start) {\textarabic{هل غُّذِيت}\\ \textarabic{البادئة}\\ \textarabic{خلال \LR{3} أيام؟}};
  \node [block, right of=decision_start, text width=7em, node distance=4cm] (feed_no_branch)
      {\textarabic{غّذِ البادئة مرتين:}\\ \textarabic{قبل \LR{\qty{48}{\hour}}}\\ \textarabic{وقبل \LR{\numrange{6}{12}~\unit{\hour}}}};
  \node [block, below of=feed_no_branch, text width=7em, node distance=2.7cm] (feed_yes_branch)
      {\textarabic{غّذِ البادئة}\\ \textarabic{بمدة \LR{\numrange{6}{12}~\unit{\hour}}}\\ \textarabic{قبل صنع العجين.}};
  \node [block, right of=feed_no_branch, text width=7em, node distance=4cm] (high_ratio)
    {\textarabic{اعتمد نسبة \LR{1:10:10}:}\\ \textarabic{\LR{10 g} بادئة،}\\ \textarabic{\LR{100 g} دقيق،}\\ \textarabic{\LR{100 g} ماء.}};
  \node [block, right of=feed_yes_branch, text width=7em, node distance=4cm] (low_ratio)
  {\textarabic{اعتمد نسبة \LR{1:5:5}:}\\ \textarabic{\LR{10 g} بادئة،}\\ \textarabic{\LR{50 g} دقيق،}\\ \textarabic{\LR{50 g} ماء.}};
  \node [decision, below of=high_ratio, node distance=6cm] (size_check)
    {\textarabic{فقاعات؟}\\ \textarabic{وزيادة في الحجم؟}};
  \node [decision, below of=decision_start, node distance=6cm] (smell_check)
    {\textarabic{رائحة خليكية}\\ \textarabic{أم لبنية زبادية؟}};
  \node [success, below of=init, node distance=6cm] (make_dough)
    {\textarabic{جهِّز العجين}};

  \path [line] (init) -- (decision_start);
  \path [line] (decision_start) -- node {\textarabic{لا}} (feed_no_branch);
  \path [line] (decision_start) -- node[below=2pt]{\textarabic{نعم}} (feed_yes_branch.north west);
  \path [line] (feed_yes_branch) -- (low_ratio);
  \path [line] (feed_no_branch) -- (high_ratio);
  \path [line] (high_ratio) -- node[anchor=east, above=2pt] {} ++(2.2,0) |-(size_check);
  \path [line] (low_ratio) -- (size_check);
  \path [line] (size_check) -- node {\textarabic{لا}} (feed_yes_branch.south east);
  \path [line] (size_check) -- node {\textarabic{نعم}} (smell_check);
  \path [line] (smell_check) -- node {\textarabic{لا}} (feed_yes_branch.south west);
  \path [line] (smell_check) -- node {\textarabic{نعم}} (make_dough);

  % braces
  \draw[BC]   (size_check.south) --
      node[below=1em]{\textarabic{تحّقق من جهوزية البادئة قبل الاستخدام}}(smell_check.south);
\end{tikzpicture}

  \caption[تحديد جهوزية بادئة الخبز الحامض]{مخطط تدفق يريك كيف تحدد جهوزية بادئة الخبز الحامض للاستخدام. احرص على الانتظار \LR{\qtyrange{6}{12}{\hour}} على الأقل بعد تغذية بادئتك قبل التحقق من جهوزيتها. ولتقييمها، انظر إلى زيادة حجم بادئتك وملمسها الإسفنجي، وسجّل رائحتها. فالعوامل الثلاثة كلها مهمة لتقييم مستوى نشاط بادئتك تقييمًا سليمًا. والبادئة النشيطة أساس مهم لتخمير عجين ناجح.}%
  \label{fig:sourdough-starter-readiness}
\end{flowchart}

والعلامة الرئيسية التي يجب مراقبتها هي الفقاعات التي تراها في برطمان بادئتك. فهي علامة على أن الخميرة تستقلب عجينك وتنتج \LR{\ch{CO2}}. وينحبس \LR{\ch{CO2}} في بنية عجينك، فيظهر على حواف الوعاء.

وانتبه أيضًا إلى زيادة حجم عجينك. فمقدار الزيادة في الحجم غير ذي أهمية. يزعم بعض الخبازين أنها تضاعف مرتين أو ثلاثًا أو أربعًا. ومقدار الزيادة يعتمد على كائناتك الدقيقة، وعلى الدقيق الذي تستخدمه لصنع البادئة. فدقيق القمح يحوي غلوتينًا أكثر، وستكون الزيادة في الحجم أكبر بسببه. وفي الوقت نفسه فالكائنات الدقيقة ليست أكثر نشاطًا على الأرجح مقارنة بوجودها في خبز حامض الجاودار. ولا يمكنك إلا أن تدّعي أن كائنات دقيق القمح الدقيقة قد تكون أفضل في تفكيك الغلوتين من كائنات الجاودار الدقيقة. وهذا من الأسباب التي دفعتني إلى تغيير دقيق بادئة الخبز الحامض عندي كثيرًا. فقد أملكت أن أنشئ بادئة شاملة تستطيع تخمير أنواع دقيق مختلفة.\footnote{سواء كان هذا ينجح، لا يمكنني الجزم علميًا. فالكائنات الدقيقة التي استقرت مرة ما تلبث قوية ولا تدع كائنات دقيقة أخرى تدخل. وقد صُنعت بادئتي أول مرة بدقيق الجاودار؛ فمن المرجح أن غالبية كائناتي الدقيقة من مصدر الجاودار.}

وأنفك أداة عظيمة أيضًا لتحديد جهوزية البادئة. وبحسب ميكروبيوم بادئتك ستنتبه إما إلى رائحة حمض اللبنيك أو حمض الخليك. فلحمض اللبنيك نفحات لبنية زبادية، ولحمض الخليك نفحات خليكية لاذعة قوية. ويصفه بعضهم برائحة الغراء أو الأسيتون. ودمج القرائن البصرية من زيادة الحجم والتجاويف مع الرائحة هو أفضل طريقة لتحديد جهوزية البادئة.

وفي حالات نادرة قد يكون دقيقك معالجًا بما يمنع نمو الكائنات الدقيقة. وقد يحدث ذلك إذا كان الدقيق غير عضوي واستخدم المزارع عوامل كيميائية حيوية كثيرة. وفي تلك الحالة جرّب ببساطة مرة أخرى بدقيق مختلف. وعشرة أيام مدة جيدة للانتظار قبل إعادة المحاولة.

ومن المنهجيات التي يستخدمها بعض الخبازين ما يسمى \textenglish{float test} (اختبار الطفو). فكرة الاختبار أن تأخذ قطعة من بادئة الخبز الحامض وتضعها فوق بعض الماء؛ فإن كان العجين مليئًا بالغاز طفا فوق الماء، وإن لم يكن جاهزًا لم يستطع الطفو وغاص إلى القاع. لكن هذا الاختبار لا يعمل مع كل أنواع الدقيق؛ فدقيق الجاودار مثلًا لا يحتفظ بالغاز كما يحتفظ به دقيق القمح، ولذا قد لا يطفو في بعض الحالات. ولهذا أنا شخصيًا لا أستخدم هذا الاختبار ولا أستطيع التوصية به.

وحين ترى بادئتك جاهزة أوصي بتغذيتها تغذية أخيرة، ثم تكون جاهزًا لصنع عجينك في المساء أو في اليوم التالي. وللتعليمات عن كيفية صنع أول عجين لك، انقل إلى الفصول التالية (\ref{ch:wheat-sourdough} و~\ref{ch:non-wheat-sourdough}) من هذا الكتاب.

وإن فشل رغيفك الأول، فالأرجح أن التخمير لم يسر كما هو متوقع. وفي كثير من الحالات يكون السبب هو بادئة الخبز الحامض. فلربما لم يتح بعد التوازن الأمثل بين البكتيريا والخميرة. وفي تلك الحالة الحل الجيد هو الاستمرار في تغذية بادئتك مرة واحدة يوميًا. فمع كل تغذية تصبح بادئتك أفضل في تخمير الدقيق، وتتأقلم الكائنات الدقيقة أكثر فأكثر مع البيئة. وينبغي أيضًا أن تنظر في فصل بادئة الخبز الحامض الصلبة \textenglish{(stiff starter)} في هذا الكتاب؛ فالبادئة الصلبة تساعد على تعزيز الجزء الخميري من بادئتك وموازنة التخمير.

\section{العناية بالبادئة}

\begin{flowchart}[!htb]
\centering
  % fig-starter-maintenance-ar.tex — Arabic edition of the sourdough
% starter maintenance flowchart (ch. 5, المخطط الانسيابي 5.3)
%
% Original: The Sourdough Framework by Hendrik Kleinwächter,
% the-sourdough-framework.com, CC BY-SA 4.0. The Arabic edition inherits
% CC BY-SA 4.0.
%
% R1-FIGURES route: re-derived from the upstream post-#492 source
% figures/fig-starter-maintenance.tex (node styles from
% figures/flowcharts_tikz.tex and the node/edge geometry verbatim), with
% the label text recovered from the accepted edition's render of p. 42
% (book-ar-fixed.pdf, b78c2a28…) — label recovery against a visual
% reference, not retranslation. Line splits inside nodes follow the
% accepted edition's render exactly (verified per node against the
% 300-dpi render and per-character coordinate dumps). Arabic label text
% is wrapped in \textarabic{}; digit/Latin runs ride \LR{} islands
% (Ruling 1 uniform-Pagella figure face, ACCEPTED 2026-11-15). The two
% feed-duration nodes render the range with a single trailing unit
% («8–12 h»), as the accepted render's per-char positions show (same
% construct as the FC 5.2 nodes).

\begin{tikzpicture}[node distance = 3cm, auto]
  \node [start] (init) {\textarabic{اصنع}\\ \textarabic{عجين الخبز}};
  \node [decision, below of=init, node distance=3.5cm] (all_starter_used) {\textarabic{هل استُخدمت}\\ \textarabic{كل البادئة؟}};
  \path [line] (init) -- (all_starter_used);
  \node [block, right of=init, node distance=3cm] (use_dough) {\textarabic{خذ \LR{10 g}}\\ \textarabic{من عجين}\\ \textarabic{الخبز}};
  \node [block, right of=all_starter_used, node distance=3cm] (use_starter) {\textarabic{خذ كل شيء}\\ \textarabic{عدا \LR{10 g}}\\ \textarabic{من بادئتك}};
  \path [line] (all_starter_used) -- node {\textarabic{نعم}} (use_dough);
  \path [line] (all_starter_used) -- node {\textarabic{لا}} (use_starter);
  \node [block, right of=use_dough, node distance=3cm] (feed_starter) {\textarabic{غّذِ باستخدام}\\ \textarabic{نسبة \LR{1:5:5}}};
  \path [line] (use_dough) -- (feed_starter);
  \path [line] (use_starter) -- (feed_starter);
  \node [decision, right of=feed_starter, node distance=3cm] (bake_next_day_check) {\textarabic{هل أخبزت}\\ \textarabic{غدًا؟}};
  \path [line] (feed_starter) -- (bake_next_day_check);
  \node [success, right of=bake_next_day_check, node distance=3.5cm]
      (make_bread_dough) {\textarabic{اصنع عجين}\\ \textarabic{الخبز مجددًا}\\ \textarabic{بعد \LR{\numrange{8}{12}~\unit{\hour}}}};
  \path [line] (bake_next_day_check) -- node {\textarabic{نعم}} (make_bread_dough);
  \node [decision, right of=use_starter, node distance=3cm] (bake_next_week_check) {\textarabic{هل ستخبز خلال}\\ \textarabic{الأسبوعين القادمين؟}};
  \node [block, right of=bake_next_week_check, node distance=3.5cm] (store_fridge) {\textarabic{خزّن البادئة}\\ \textarabic{في الثلاجة}\\ \textarabic{عند \LR{\qty{4}{\degreeCelsius} (\qty{40}{\degF})}}};
  \path [line] (bake_next_week_check) -- node {\textarabic{نعم}} (store_fridge);
  \node [success, right of=store_fridge, node distance=3cm] (feed_after_fridge) {\textarabic{غّذِ مجددًا}\\ \textarabic{بنسبة \LR{1:5:5}}\\ \textarabic{قبل صنع العجين}\\ \textarabic{بمدة \LR{\numrange{8}{12}~\unit{\hour}}}};
  \path [line] (store_fridge) -- (feed_after_fridge);
  \path [line] (bake_next_day_check) -- node {\textarabic{لا}} (bake_next_week_check);
  \node [decision, below of=use_starter, node distance=3cm] (freezer_check) {\textarabic{هل لديك}\\ \textarabic{مُجمِّد؟}};
  \path [line] (bake_next_week_check) -- node {\textarabic{لا}} (freezer_check);
  \node [block, right of=freezer_check, node distance=3cm] (dry_starter) {\textarabic{جفّف}\\ \textarabic{بادئتك}};
  \node [block, below of=dry_starter, node distance=3cm] (freeze_starter) {\textarabic{جمّد}\\ \textarabic{بادئتك}};
  \path [line] (freezer_check) -- node {\textarabic{لا}} (dry_starter);
  \path [line] (freezer_check) -- node {\textarabic{نعم}} (freeze_starter);
  \node [success, right of=dry_starter, node distance=3.5cm] (reactivate_freezer) {\textarabic{أعد تنشيط البادئة}\\ \textarabic{خلال \LR{3} أيام}\\ \textarabic{بتغذيات يومية}\\ \textarabic{بنسبة \LR{1:5:5}}};
  \path [line] (dry_starter) -- (reactivate_freezer);
  \path [line] (freeze_starter) -- (reactivate_freezer);
\end{tikzpicture}

  \caption[مخطط عناية بادئة الخبز الحامض الانسيابي]{مخطط تدفق كامل يريك كيف تُجري عناية سليمة ببادئة الخبز الحامض. يمكنك استخدام قطعة من عجينك كبادئة تالية. ويمكنك أيضًا استخدام البادئة المتبقية وتغذيتها مرة أخرى. واختر الخيار الأنسب لجدولك الخاص. ويفترض المخطط أنك تستخدم بادئة عند مستوى ترطيب \LR{\qty{100}{\percent}}. فاضبط محتوى الماء وفقًا لذلك عند استخدام بادئة صلبة.}%
  \label{fig:sourdough-maintenance-process}
\end{flowchart}

لقد صنعت بادئة الخبز الحامض ورغيفك الأول. فكيف تُجري العناية ببادئتك؟ هناك طرق عناية لا تُحصى. ويهوّم بعض الناس ببادئتهم حتى يغذّونها يوميًا. والمفتاح لفهم كيفية إجراء العناية السليمة هو فهم ما يحدث لبادئتك بعد استخدامها لصنع عجين. فأي بادئة متبقية لديك، أو قطعة صغيرة من عجين الخبز، يمكن أن تخدم لصنع بادئتك التالية.\footnote{أنا كثيرًا ما أستخدم كل بادئتي لصنع عجين. فإن طلبت الوصفة \LR{\qty{50}{\gram}} من البادئة صنعت مسبقًا \LR{\qty{50}{\gram}} بادئة بالضبط؛ أي لا يتبقى لديّ بادئة. وحينها آخذ قطعة صغيرة من العجين في نهاية فترة التخمير وأستخدمها لإعادة تنمية بادئتي.}

وكما شرحنا سابقًا، بادئتك متأقلمة مع تخمير الدقيق. والكائنات الدقيقة فيها صلبة البأس جدًا؛ فهي تحجب المسبّبات الخارجية وسائر الكائنات الدقيقة. ولهذا، عند شراء بادئة خبز حامض ستحافظ على الكائنات الدقيقة الأصلية. فمن المرجح أنها لن تتغير في بادئتك؛ فهي تتفوق على سائر الكائنات الدقيقة في تخمير الدقيق. وكل شيء في الطبيعة يبدأ عادة بالتحلل بعد حين. غير أن كائنات بادئتك الدقيقة تملك آليات دفاع قوية جدًا. وفي نهاية المطاف، يمكن مقارنة بادئة الخبز الحامض بالمخللات. فالمخللات ثبت أنها تبقى صالحة مدة طويلة جدًا~\cite{pickled+foods+expiration}. فحمضية بادئة الخبز الحامض سامة للغاية للكائنات الدقيقة الأخرى. غير أن الخمائر والبكتيريا تأقلمت مع العيش في البيئة عالية الحمضية. قارن ذلك بمعدتك، حيث تحيد الحمضية كثيرًا من المسبّبات الممكنة. وما دامت بادئتك تملك غذاءً كافيًا فستتفوق على سائر الكائنات الدقيقة. وحين ينفد غذاء البادئة تبدأ الكائنات الدقيقة بالبوغ. فهي تستعد لفترة انعدام غذاء، ثم تعاود النشاط لحظة حضور غذاء جديد. والأبواغ صلبة المقاومة جدًا ويمكنها البقاء في ظروف قاسية. وزعم العلماء أنهم وجدوا أبواغًا عمرها 250 مليون سنة ما تزال نشيطة~\cite{old+spores}. ومع ذلك فهي في صورة الأبواغ أهش أمام المسبّبات الخارجية، كالعفن. ومع ذلك، في الظروف المثالية يمكن للأبواغ أن تبقى مدة طويلة.

لكن ما دامت في بيئة بادئتك فهي تعيش في بيئة محمية جدًا. ويصعب على الفطريات والبكتيريا الأخرى أن تحلل كتلة بادئتك المتبقية. ولقد رأيت حالات نادرة جدًا ماتت فيها البادئة فعلًا. فمن المستحيل تقريبًا قتل بادئة.

ما يحدث مع ذلك هو أن توازن الخمائر والبكتيريا يتغير في بادئتك. فالبكتيريا أقدر على العيش في بيئة حمضية. وهذه مشكلة حين تصنع عجينًا آخر؛ إذ تريد التوازن السليم بين الرغوة والنفحات الحامضة. وحين سبَت البادئة مدة طويلة فمن المرجح ألا يكون لديك توازن مرغوب من الكائنات الدقيقة. وعلاوة على ذلك، بحسب مدة سبات بادئتك قد لا يتبقى لديك إلا كائنات دقيقة بَوَّغت. فهذه التغذيات القليلة ستساعد في إعادة بادئة الخبز الحامض إلى شكلها السليم.

فيما يلي بعض السيناريوهات التي ستساعدك على إجراء عناية سليمة بالبادئة، بحسب موعد خبزك القادم.

\begin{description}
\item[أريد أن أخبز مرة أخرى في اليوم التالي:]
خذ ببساطة ما تبقى لديك من بادئة وغذِّها مرة أخرى. وإن استنفدت كل بادئتك يمكنك قطع قطعة من عجينك؛ فالعجين نفسه ليس سوى بادئة عملاقة. وأوصي بالنسبة \LR{1:5:5} كما ذكرنا سابقًا. فخذ جزءًا واحدًا من البادئة، وغذِّها بخمسة أجزاء من الدقيق وخمسة أجزاء من الماء. وإن كان الجو حارًا جدًا حيث تسكن، أو إن أردت صنع الخبز نحو 24 ساعة بعد آخر تغذية، فغيّر النسبة. في تلك الحالة سأعتمد نسبة \LR{1:10:10}. وأحيانًا لا توجد لديّ بادئة كافية؛ حينها أستخدم حتى نسبة \LR{1:50:50} أو \LR{1:100:100}. وبحسب كمية الدقيق الجديد التي تغذّيها يستغرق الأمر مدة أطول حتى تصبح بادئتك جاهزة مرة أخرى.

\item[أريد أن آخذ استراحة وأخبز الأسبوع القادم:]
خذ بادئتك المتبقية وضعها في ثلاجتك؛ فستبقى صالحة مدة طويلة جدًا. والشيء الوحيد الذي أراه يحدث هو جفاف السطح في الثلاجة. لذا أوصي بغمر البادئة في قليل من الماء؛ فهذه الطبقة المائية الإضافية توفر حماية جيدة من جفاف الجزء العلوي. فالعفن هوائي التنفس، فلا يستطيع النمو بكفاءة تحت الماء~\cite{mold+anaerobic}. وقبل استخدام البادئة مرة أخرى، امزج السائل في العجين أو صفِّه ببساطة. وإن صفّيت السائل يمكنك استخدامه لصنع صلصة حارة مخمَّرة بالحمض اللبنيكي مثلًا.

كلما كان الجو أبرد بقيت مدة أطول على توازن جيد بين الخمائر والبكتيريا. وبصورة عامة، كلما ارتفعت الحرارة كان التخمير أسرع، وكلما انخفضت صارت العملية كلها أبطأ. وفأقل من \LR{\qty{4}{\degreeCelsius}} يتوقف تخمير البادئة شبه تام. وتعتمد سرعة التخمير في درجات الحرارة المنخفضة على سلالات الخمائر والبكتيريا البرية التي نمّيتها.

\item[أريد أن آخذ استراحة عدة أشهر:]
قد يكون تجفيف بادئتك هو الخيار الأفضل لحفظها في هذه الحالة. فبما أنك تزيل الرطوبة والغذاء، ستبَوِّغ كائناتك الدقيقة. وبما أنه لا رطوبة، تصمد الأبواغ أمام المسبّبات الآخرين بثبات كبير. والبادئة المجففة قد تبقى صالحة سنوات.

خذ بادئتك ببساطة واخلطها بالدقيق. وحاول تفتيت البادئة قدر الإمكان. وأضف مزيدًا من الدقيق باستمرار حتى تلاحظ عدم بقاء أي رطوبة. وضع بادئة الدقيق في مكان جاف في بيتك ودَعها تجف أكثر. وإن كان لديك مجفف كهربائي يمكنك استخدامه لتسريع العملية. اضبطه على نحو \LR{\qty{30}{\degreeCelsius}} وجفّف البادئة \LR{12--20} ساعة. في اليوم التالي تكون بادئتك قد جفت قليلًا. وهي في حالة هشة إذ ما زالت فيها رطوبة. أضف مزيدًا من الدقيق لتسريع التجفيف، وكرر ذلك لمدة \num{2} أيام أخرى حتى تشعر بعدم بقاء أي رطوبة. وهذا مهم وإلا بدأت بالعفن. وبمجرد الانتهاء خزّن البادئة في وعاء محكم الإغلاق. أو يمكنك المضي وتجميد البادئة المجففة. وكلا الخيارين يعمل بشكل ممتاز. وبادئتك المبوَّغة الآن تنتظر تغذيتك التالية. وإن توفرت أضف بعض أكياس السيليكا إلى الوعاء لامتصاص الرطوبة الزائدة أكثر.

في البداية، كانت تحتاج بادئتي إلى نحو ثلاثة أيام لتعود حية بعد تجفيفها وإعادة تنشيطها. وإن فعلت الشيء نفسه اليوم فأحيانًا تكون بادئتي جاهزة بعد تغذية واحدة. ويبدو أن الكائنات الدقيقة تتأقلم؛ فالناجية من هذا الصدمة تصبح مهيمنة فيما بعد.
\end{description}

خلاصة القول، إن طريقة العناية التي تختارها تعتمد على موعد خبزك القادم. والهدف من كل تغذية جديدة هو ضمان أن بادئتك تملك التوازن المرغوب بين الخمائر والبكتيريا عند صنع العجين. ولا داعي لتغذية بادئتك يوميًا، إلا إذا لم تنضج بعد. وفي تلك الحالة، ستساعد كل تغذية تالية في جعل بادئتك أكثر برعة في تخمير الدقيق.
